\section{Subagent 系统设计}

\subsection{为什么需要 Subagent？}

\begin{importantbox}{三大动机}
\begin{enumerate}
    \item \textbf{隔离上下文}：避免所有信息塞入单个 Master Agent 的 Context Window 导致溢出。
    \item \textbf{并发执行}：多个独立任务可以由不同 Subagent 同时处理。
    \item \textbf{更长的交付物}：Master Agent 将繁重工作委派给 Subagent，自己只做最终的 Reduce/汇总。
\end{enumerate}
\end{importantbox}

\subsection{五个 Subagent 操作}

Codex 的 Subagent 系统由五个 Function Tool 组成：

\begin{enumerate}
    \item \textbf{Spawn Agents}：创建子 Agent。
    \item \textbf{Send Inputs}：向已有子 Agent 发送新指令。
    \item \textbf{Wait}：阻塞等待一个或多个子 Agent 达到最终状态。
    \item \textbf{Close Agents}：显式关闭/下档子 Agent。
    \item \textbf{Resume}：恢复一个已关闭或退出的子 Agent，沿其原有上下文继续工作。
\end{enumerate}

\begin{warningbox}{Subagent 的通信方向是单向的}
子 Agent 只在\textbf{结束时}才产生输出。当前的结构设计中，子 Agent 没有主动与主 Agent 交互的模式——所有通信都是主 Agent 向子 Agent 发送。
\end{warningbox}

\subsection{Agent Nickname 机制}

Codex Runtime 维护了一个 Agent Nicknames 池，主要是历史科学家和学者的名字。创建子 Agent 时：
\begin{enumerate}
    \item 优先查看配置文件中是否自定义了 nickname。
    \item 否则从内置名字池中选取一个未被占用的名字。
    \item Nickname 用于前端展示，所有内部交互使用 Agent ID。
\end{enumerate}

\subsection{Spawn Agent 的返回值与上下文感知}

\begin{knowledgebox}{模型如何感知创建的 Subagent？}
Function Calling 的 Tool 定义中没有返回值描述。那么模型如何知道自己创建了哪个 Subagent？

答案在于：\textbf{Spawn Agent 的返回值}（包含 agent\_id 和 nickname）作为 Function Call Output 被 Runtime 回填到模型上下文中。模型在后续推理时自然能从上下文中读取到这些信息，从而建立起"我创建了哪个 Agent、它的 ID 是什么"的关联。
\end{knowledgebox}

\subsection{Wait Agent 的语义}

\begin{itemize}
    \item Wait 不是 \texttt{sleep}，而是对 Agent 状态流的\textbf{订阅}——更像是 \texttt{join} / \texttt{await} 操作。
    \item 它会阻塞 Master Agent 的后续推理，直到目标 Subagent 达到最终状态。
    \item 有超时机制（最长约一小时）。
    \item \textbf{何时使用}：仅当 Master Agent 的下一步必须依赖 Subagent 的结果时才触发。如果自己还有别的事情可做，不要急于 Wait。
\end{itemize}

\begin{importantbox}{Spawn Agent 的触发条件}
Codex 对 Spawn Agent 的使用设置了\textbf{两条强指令}：
\begin{enumerate}
    \item 只有用户\textbf{显式要求} Subagent 分派或并发执行时，才创建 Subagent。
    \item 例外：如果是对 Codebase 做\textbf{深度、全面的 Research 或系统分析}，不需要用户许可即可创建 Subagent。
\end{enumerate}
因此，想测试 Subagent 流程，需要在 Query 中显式提到"用 subagent 并发完成"。
\end{importantbox}

\subsection{Subagent 完成后的通知机制}

如果 Master Agent 没有调用 Wait，Subagent 完成后并不会"失联"。Runtime 会在后台自动将 Subagent 的完成消息\textbf{以 User 身份}注入到 Master Agent 的上下文中，确保信息不丢失。

\subsection{本章小结}

Subagent 系统通过上下文隔离和并发执行，突破了单一 Context Window 的限制。五个操作原语（Spawn / Send Inputs / Wait / Close / Resume）构成了完整的子 Agent 生命周期管理。


